\documentclass[a4paper]{article}

\usepackage{graphicx}
\usepackage{amsmath,amssymb}
\usepackage{minted}
\usepackage{multirow}
\usepackage{float}
\usepackage{todonotes}
\usepackage{forest}
\usepackage{xstring}
\usepackage{xcolor}
\usepackage[most]{tcolorbox}
\usepackage{xparse}
\usepackage{natbib}

\usetikzlibrary{graphs,graphdrawing}
\usegdlibrary{layered}

% for external graphviz
\input{/home/donkarlo/Dropbox/repo/nd_ai_project/src/nd_ai/knoweldge_representation/language/formal/computer/programming/tex/latex/code_clips/external_graphviz.tex}

% for tags
\input{/home/donkarlo/Dropbox/repo/nd_ai_project/src/nd_ai/knoweldge_representation/language/formal/computer/programming/tex/latex/parts/control_sequence/kind/semtag/semtag.tex}

% for links
\input{/home/donkarlo/Dropbox/repo/nd_ai_project/src/nd_ai/knoweldge_representation/language/formal/computer/programming/tex/latex/code_clips/seehere.tex}

\title{Memory trace or engram}
\date{}
\author{Mohammad Rahmani}

\begin{document}
    \maketitle
    An engram is a unit of cognitive information imprinted in a physical substance, theorized to be the means by which memories are stored[1] as biophysical or biochemical[2] changes in the brain or other biological tissue, in response to external stimuli \seeherefooter{https://en.wikipedia.org/wiki/Engram_(neuropsychology)}.
    
    \seeherefooter{https://pubmed.ncbi.nlm.nih.gov/31896692/}
    
    \seeherefooter{file:///home/donkarlo/Dropbox/repo/nd_shared_research/bibliography/corpora/1756-6606-5-32.pdf}
    
    \cite{sakaguchi-2012-catching-the-engram-strategies-to-examine-the-memory-trace}
    
    \bibliographystyle{plainnat}
    \bibliography{/home/donkarlo/Dropbox/repo/nd_shared_research/bibliography/bibliography.bib}
\end{document}